\documentclass[11pt]{article}
\usepackage[margin=0.68in]{geometry}
\usepackage{lmodern}
\usepackage{microtype}
\usepackage{enumitem}
\usepackage[table]{xcolor}
\usepackage{titlesec}
\setlength{\parindent}{0pt}
\setlength{\parskip}{0.24em}
\setlist[itemize]{leftmargin=1.3em,itemsep=0.05em,topsep=0.08em}
\titleformat{\section}{\large\bfseries}{}{0pt}{}
\titlespacing*{\section}{0pt}{0.5em}{0.18em}
\pagestyle{plain}
\newenvironment{playpassage}{\begin{quote}\small\setlength{\parskip}{0.04em}\setlength{\parindent}{0pt}}{\end{quote}}
\newcommand{\speaker}[1]{\textbf{#1:}}
\begin{document}
\begin{center}
{\Large\bfseries U1L17SP --- Text Supplement}\par
{\LARGE\bfseries Two Private Arguments}\par
{\large Close-reading anchors from \textit{Macbeth} and \textit{Julius Caesar}}
\end{center}
\fcolorbox{black}{gray!8}{\begin{minipage}{0.94\linewidth}\small
\textbf{How to use this supplement}

You may use \textbf{anything accurate from either play}. The passages below are lines worth focusing on, not limits on your evidence. Quote exact words when they are printed here. You may accurately paraphrase other events you remember from watching the plays. Do not invent quotations from memory.
\end{minipage}}

\section*{Passage 1: Macbeth, Act 1, Scene 7 --- Macbeth Before the Murder}
\textbf{Context:} Duncan is Macbeth's king, kinsman, and guest. Alone, Macbeth considers consequences, duties, Duncan's character, and what drives his own intent before Lady Macbeth returns.
\begin{playpassage}
\speaker{Macbeth} If it were done when 'tis done, then 'twere well\par
It were done quickly: if the assassination\par
Could trammel up the consequence, and catch\par
With his surcease success; that but this blow\par
Might be the be-all and the end-all here,\par
But here, upon this bank and shoal of time,\par
We'ld jump the life to come. But in these cases\par
We still have judgment here; that we but teach\par
Bloody instructions, which, being taught, return\par
To plague the inventor.\par

He's here in double trust;\par
First, as I am his kinsman and his subject,\par
Strong both against the deed; then, as his host,\par
Who should against his murderer shut the door,\par
Not bear the knife myself.\par

Besides, this Duncan\par
Hath borne his faculties so meek, hath been\par
So clear in his great office, that his virtues\par
Will plead like angels, trumpet-tongued, against\par
The deep damnation of his taking-off.\par

I have no spur\par
To prick the sides of my intent, but only\par
Vaulting ambition, which o'erleaps itself\par
And falls on the other.
\end{playpassage}
\section*{Words in context}
\begin{itemize}
\item \textbf{trammel up}: contain; \textbf{surcease}: death; \textbf{jump}: risk;
\item \textbf{judgment here}: consequences or justice in this life; \textbf{double trust}: multiple obligations;
\item \textbf{faculties}: powers and duties of office; \textbf{taking-off}: killing; \textbf{spur}: motive that drives action.
\end{itemize}

\newpage
\section*{Passage 2: Julius Caesar, Act 2, Scene 1 --- Brutus in the Orchard}
\textbf{Context:} Alone before joining the conspiracy, Brutus considers Caesar's possible future power. He says he has no personal cause against Caesar and has not known Caesar's passions to overpower his reason.
\begin{playpassage}
\speaker{Brutus} It must be by his death: and for my part,\par
I know no personal cause to spurn at him,\par
But for the general. He would be crown'd:\par
How that might change his nature, there's the question.\par
It is the bright day that brings forth the adder;\par
And that craves wary walking. Crown him? That;\par
And then, I grant, we put a sting in him,\par
That at his will he may do danger with.\par
The abuse of greatness is, when it disjoins\par
Remorse from power: and, to speak truth of Caesar,\par
I have not known when his affections sway'd\par
More than his reason. But 'tis a common proof,\par
That lowliness is young ambition's ladder,\par
Whereto the climber-upward turns his face;\par
But when he once attains the upmost round,\par
He then unto the ladder turns his back.\par

So Caesar may.\par
Then, lest he may, prevent. And, since the quarrel\par
Will bear no colour for the thing he is,\par
Fashion it thus; that what he is, augmented,\par
Would run to these and these extremities:\par
And therefore think him as a serpent's egg\par
Which, hatch'd, would, as his kind, grow mischievous,\par
And kill him in the shell.
\end{playpassage}
\section*{Words in context}
\begin{itemize}
\item \textbf{for the general}: for the public good; \textbf{remorse}: pity or restraint;
\item \textbf{affections}: passions or desires; \textbf{common proof}: common experience;
\item \textbf{prevent}: act beforehand; \textbf{colour}: plausible justification;
\item \textbf{augmented}: increased in power; \textbf{mischievous}: harmful.
\end{itemize}

\section*{Broader evidence and fair comparison}
You may test either private argument against later conduct and consequences. For Macbeth, possible reminders include Lady Macbeth's pressure, Duncan's murder, further violence, the apparitions, and Macbeth's final choices. For Brutus, possibilities include the forged letters, the assassination, his funeral defense, the civil war, and his later judgments. These are optional reminders, not diagnoses.

A character may admit some truth and still reason dishonestly elsewhere. A mistaken prediction is not automatically self-deception. Explain what each man knows, wants, assumes, and does with contrary pressure before comparing them.
\end{document}
